\documentclass[agenda]{sfs-2487-2024}

% Kokouskutsu esityslistoineen. Vaihtoehto [agenda] numeroi otsikot
% loppupisteellisinä ("1. Otsikko"), mikä on standardin 6.4.1 salliman
% vakiintuneen tavan mukaista esityslistoissa ja pöytäkirjoissa.

\logo{\includegraphics{logo-organisaatio}}
\doctype{Kokouskutsu}
\date{6.5.2024}
\author{Virve Virtanen}

\subject{Digiprojekti}
\title{Projektiryhmän kokous 5/2024}

\contactinfo{Organisaatio Oy}{%
  Katuosoite\\
  Postinumero Postitoimipaikka\\
  Puhelinnumero\\
  www-osoite}

\begin{document}

\maketitle

Projektiryhmän seuraava kokous pidetään maanantaina 13.5.2024
klo 12.30--13.45 verkkokokouksena. Kokouslinkki lähetetään
osallistujille kalenterikutsussa.

\section*{Esityslista}

\section{Kokouksen avaus}

\section{Asiakaspalautteen läpikäynti ja tehtävät muutokset}

Yhteenveto asiakaspalautteesta on kutsun liitteenä.

\section{Etusivun uudistaminen}

Käsitellään viestintäosaston sihteerin ehdotus etusivulle
tehtävistä muutoksista.

\section{Seuraavat kokoukset}

\attachments{Yhteenveto asiakaspalautteesta}

\distribution{Projektiryhmän jäsenet}

\forinformation{Johtoryhmä}

\makecontactinfo

\end{document}
